\documentclass[10pt,a4paper,landscape]{article}
\usepackage[utf8]{inputenc}
\usepackage[french]{babel}
\usepackage[T1]{fontenc}
\usepackage{multicol}
\usepackage{amsmath,amssymb}
\usepackage{xcolor}
\usepackage{tcolorbox}
\usepackage{geometry}
\usepackage{enumitem}
\geometry{margin=1cm}

\setlist{nosep,leftmargin=*}

% Couleurs
\definecolor{defvert}{RGB}{0,153,76}
\definecolor{exempleorange}{RGB}{255,102,0}
\definecolor{importantrouge}{RGB}{204,0,0}
\definecolor{titrebleu}{RGB}{0,102,204}

% Boîtes compactes
\newtcolorbox{defbox}{
    colback=green!5,
    colframe=defvert,
    fonttitle=\bfseries\small,
    title=�� DÉFINITION,
    boxrule=0.5mm,
    left=2mm,right=2mm,top=1mm,bottom=1mm
}

\newtcolorbox{algobox}{
    colback=orange!5,
    colframe=exempleorange,
    fonttitle=\bfseries\small,
    title=⚙️ ALGORITHME,
    boxrule=0.5mm,
    left=2mm,right=2mm,top=1mm,bottom=1mm
}

\newtcolorbox{importbox}{
    colback=red!5,
    colframe=importantrouge,
    fonttitle=\bfseries\small,
    title=⚠️ IMPORTANT,
    boxrule=0.5mm,
    left=2mm,right=2mm,top=1mm,bottom=1mm
}

\newtcolorbox{complexbox}{
    colback=blue!5,
    colframe=titrebleu,
    fonttitle=\bfseries\small,
    title=⏱️ COMPLEXITÉS,
    boxrule=0.5mm,
    left=2mm,right=2mm,top=1mm,bottom=1mm
}

\pagestyle{empty}

\begin{document}

\begin{center}
\LARGE\textbf{\color{titrebleu}FICHE DE RÉVISION - À SAVOIR / FAIRE}\\
\large\textit{Niveaux de Langages + Types Abstraits de Données}
\end{center}

\begin{multicols}{3}
\small

\section*{\color{titrebleu}1. NIVEAUX DE LANGAGES}

\subsection*{MÉMOIRE ET VARIABLES}

\begin{defbox}
\textbf{Mémoire}
\begin{itemize}
\item \textbf{Bits} : contacts magnétiques (0 ou 1)
\item \textbf{Mots} : groupes de bits adressables
\item Exemple : 65 536 mots × 32 bits
\item \textbf{IP/PC} : compteur ordinal, pointe vers prochaine instruction
\end{itemize}
\end{defbox}

\begin{defbox}
\textbf{Cellule mémoire}
\begin{itemize}
\item Statut : allouée / non allouée
\item Contenu : valeur définie / indéfinie (?)
\end{itemize}
\end{defbox}

\begin{defbox}
\textbf{Variable}
\begin{itemize}
\item Objet contenant une valeur
\item Inspectable et modifiable
\item \textbf{Simple} : valeur primitive
\item \textbf{Composée} : record, struct
\end{itemize}
\end{defbox}

\subsection*{SYNTAXE ET SÉMANTIQUE}

\begin{importbox}
\textbf{Syntaxe} : forme du programme\\
\textbf{Sémantique} : sens, comportement
\end{importbox}

\subsection*{LANGAGE MACHINE}

\begin{defbox}
\textbf{Langage machine}
\begin{itemize}
\item Suite de bits
\item Seul langage du processeur
\item Directement exécutable
\item Illisible pour humain
\item Spécifique à l'architecture
\end{itemize}
\end{defbox}

\subsection*{LANGAGE D'ASSEMBLAGE}

\begin{defbox}
\textbf{Assemblage}
\begin{itemize}
\item Représentation lisible du machine
\item Correspondance 1:1 avec machine
\item Mnémoniques (ADD, MOV, PUSH)
\end{itemize}
\end{defbox}

\begin{importbox}
\textbf{Assembleur} : programme qui traduit assemblage → machine
\end{importbox}

\subsection*{LANGAGE DE COMPILATION}

\begin{defbox}
\textbf{Langage de compilation}
\begin{itemize}
\item Haut niveau d'abstraction
\item Universel, naturel, implémentable
\item Modèles récurrents
\item Unités indépendantes
\end{itemize}
\end{defbox}

\subsection*{COMPILATEUR}

\begin{importbox}
\textbf{Compilateur}
\begin{itemize}
\item Transforme code source → code objet
\item Années 50
\item Source (haut niveau) → Objet (bas niveau)
\end{itemize}
\end{importbox}

\begin{algobox}
\textbf{Processus}\\
Source → Compilateur → Assembleur → Assembleur → Machine
\end{algobox}

\begin{defbox}
\textbf{Bootstrapping}
\begin{enumerate}
\item Compilateur minimal en assembleur
\item Compiler version complète
\item Répéter (amélioration)
\end{enumerate}
\end{defbox}

\subsection*{INTERPRÉTEUR}

\begin{defbox}
\textbf{Interpréteur}
\begin{itemize}
\item Analyse + traduit + exécute
\item Ligne par ligne
\item Pas de code intermédiaire
\end{itemize}
\end{defbox}

\begin{importbox}
\textbf{Cycle}
\begin{enumerate}
\item Lire et analyser
\item Si correct, exécuter
\item Passer à suivante
\end{enumerate}
\end{importbox}

\subsection*{LDA (Pseudo-code)}

\begin{defbox}
\textbf{LDA}
\begin{itemize}
\item Langage de Description d'Algorithme
\item Routines de calcul
\item Gestion mémoire abstraite
\item Paradigme objet
\item Distingue fonction/procédure
\end{itemize}
\end{defbox}

\begin{importbox}
\textbf{Fonction} : retourne valeur\\
\textbf{Procédure} : effet de bord
\end{importbox}

\columnbreak

\section*{\color{titrebleu}2. STRUCTURES DE DONNÉES}

\subsection*{PILE (Stack)}

\begin{defbox}
\textbf{Pile - LIFO}
\begin{itemize}
\item Last In, First Out
\item Dernier entré = premier sorti
\end{itemize}
\end{defbox}

\begin{algobox}
\textbf{Opérations}
\begin{itemize}
\item \texttt{empiler(x)} / push
\item \texttt{dépiler()} / pop
\item \texttt{pileVide?()}
\item \texttt{sommet()} / peek
\end{itemize}
\end{algobox}

\begin{complexbox}
\textbf{Toutes} : $O(1)$
\end{complexbox}

\begin{algobox}
\small
\textbf{Empiler}
\begin{verbatim}
Si taille >= max Alors
    Erreur
Sinon
    taille++
    tab[taille] <- x
\end{verbatim}
\end{algobox}

\subsection*{FILE (Queue)}

\begin{defbox}
\textbf{File - FIFO}
\begin{itemize}
\item First In, First Out
\item Premier entré = premier sorti
\end{itemize}
\end{defbox}

\begin{algobox}
\textbf{Opérations}
\begin{itemize}
\item \texttt{enfiler(x)} / enqueue
\item \texttt{défiler()} / dequeue
\item \texttt{fileVide?()}
\end{itemize}
\end{algobox}

\begin{complexbox}
\textbf{Toutes} : $O(1)$
\end{complexbox}

\begin{importbox}
\textbf{Implémentations}
\begin{itemize}
\item Tableau circulaire
\item Double pile
\item Liste chaînée
\end{itemize}
\end{importbox}

\subsection*{ABR (Arbre Binaire Recherche)}

\begin{defbox}
\textbf{ABR}
\begin{itemize}
\item Sous-arbre gauche $\leq$ nœud
\item Sous-arbre droit $>$ nœud
\item Structure récursive
\end{itemize}
\end{defbox}

\begin{algobox}
\textbf{Accesseurs}
\begin{itemize}
\item \texttt{clé}
\item \texttt{sAG}, \texttt{sAD}
\item \texttt{père}
\end{itemize}
\end{algobox}

\begin{complexbox}
\begin{tabular}{lcc}
\textbf{Op.} & \textbf{Moy.} & \textbf{Pire} \\
\hline
Recherche & $O(\log n)$ & $O(n)$ \\
Insertion & $O(\log n)$ & $O(n)$ \\
Suppression & $O(\log n)$ & $O(n)$ \\
Min/Max & $O(\log n)$ & $O(n)$ \\
\end{tabular}
\end{complexbox}

\begin{algobox}
\small
\textbf{Rechercher}
\begin{verbatim}
Si arbre = NIL OU clé = x
    Retourner arbre
Si x < clé
    Retourner rechercher(sAG, x)
Sinon
    Retourner rechercher(sAD, x)
\end{verbatim}
\end{algobox}

\subsection*{TAS (Heap)}

\begin{defbox}
\textbf{Tas}
\begin{itemize}
\item Arbre binaire presque complet
\item \textbf{Max} : père $\geq$ fils
\item \textbf{Min} : père $\leq$ fils
\item Implémentation : tableau
\end{itemize}
\end{defbox}

\begin{algobox}
\textbf{Relations}
\begin{itemize}
\item \texttt{père(i) = i DIV 2}
\item \texttt{filsGauche(i) = 2i}
\item \texttt{filsDroit(i) = 2i+1}
\item \texttt{maximum() = tab[1]}
\end{itemize}
\end{algobox}

\begin{complexbox}
\begin{tabular}{lc}
\textbf{Opération} & \textbf{Complexité} \\
\hline
Insertion & $O(\log n)$ \\
Extraction max & $O(\log n)$ \\
Recherche max & $O(1)$ \\
Construction & $O(n)$ \\
\end{tabular}
\end{complexbox}

\columnbreak

\section*{\color{titrebleu}3. TYPES ABSTRAITS}

\subsection*{TAD - Définition}

\begin{defbox}
\textbf{TAD}
\begin{itemize}
\item Ensemble d'objets similaires
\item Ensemble d'opérations
\item Spécification $\neq$ Implémentation
\end{itemize}
\end{defbox}

\begin{importbox}
\textbf{Contraintes}
\begin{itemize}
\item Min espace mémoire
\item Min instructions
\end{itemize}
\end{importbox}

\subsection*{DICTIONNAIRE}

\begin{defbox}
\textbf{Dictionnaire}
\begin{itemize}
\item Ensemble fini ordonné
\item Opérations : insert, search, delete
\end{itemize}
\end{defbox}

\begin{complexbox}
\small
\begin{tabular}{lccc}
\textbf{Structure} & \textbf{R} & \textbf{I} & \textbf{S} \\
\hline
Tab non ord. & $O(n)$ & $O(1)$ & $O(1)$ \\
Liste non ord. & $O(n)$ & $O(1)$ & $O(1)$ \\
Tab ordonné & $O(\log n)$ & $O(n)$ & $O(n)$ \\
ABR & $O(h)$ & $O(h)$ & $O(h)$ \\
Tas & $O(n)$ & $O(\log n)$ & $O(\log n)$ \\
\end{tabular}
\end{complexbox}

\subsection*{ENSEMBLE DYNAMIQUE}

\begin{defbox}
\textbf{Ensemble Dynamique}
\begin{itemize}
\item Extension du dictionnaire
\item + min, max, pred, succ
\end{itemize}
\end{defbox}

\begin{algobox}
\textbf{Opérations}
\begin{itemize}
\item \texttt{inserer(x)}
\item \texttt{supprimer(x)}
\item \texttt{rechercher(x)}
\item \texttt{minimum()}
\item \texttt{maximum()}
\item \texttt{predecesseur(x)}
\item \texttt{successeur(x)}
\end{itemize}
\end{algobox}

\begin{importbox}
\textbf{Implémentation}\\
Liste doublement chaînée ordonnée
\begin{itemize}
\item \texttt{premier}, \texttt{dernier}
\item \texttt{précédent}, \texttt{suivant}
\end{itemize}
\end{importbox}

\subsection*{FILE DE PRIORITÉ}

\begin{defbox}
\textbf{File de Priorité}
\begin{itemize}
\item Éléments avec clé
\item Ordre de priorité sur clés
\end{itemize}
\end{defbox}

\begin{algobox}
\textbf{Opérations}
\begin{itemize}
\item \texttt{inserer(x, clé)}
\item \texttt{maximum()}
\item \texttt{extraireMax()}
\end{itemize}
\end{algobox}

\begin{complexbox}
\small
\begin{tabular}{lccc}
\textbf{Struct.} & \textbf{Ins} & \textbf{Max} & \textbf{Extr} \\
\hline
Tab non ord. & $O(1)$ & $O(n)$ & $O(n)$ \\
Liste ord. & $O(n)$ & $O(1)$ & $O(1)$ \\
\textbf{TAS} & $\mathbf{O(\log n)}$ & $\mathbf{O(1)}$ & $\mathbf{O(\log n)}$ \\
\end{tabular}
\end{complexbox}

\begin{importbox}
\textbf{Tas = structure optimale !}
\end{importbox}

\subsection*{FAMILLE D'ENSEMBLES}

\begin{defbox}
\textbf{Famille d'Ensembles}
\begin{itemize}
\item Ensemble de sous-ensembles
\item Relation d'ordre $<_X$
\end{itemize}
\end{defbox}

\begin{algobox}
\textbf{Opérations}
\begin{itemize}
\item \texttt{inserer(ensemble)}
\item \texttt{supprimer(ensemble)}
\item \texttt{appartenance(ensemble)}
\end{itemize}
\end{algobox}

\begin{importbox}
\textbf{Implémentation}\\
Arbre lexicographique (Trie)
\begin{itemize}
\item Arêtes étiquetées
\item Nœuds terminaux
\item Ordre $<_X$
\end{itemize}
\end{importbox}

\begin{complexbox}
Pour ensemble de taille $k$ :\\
\textbf{Toutes opérations} : $O(k)$
\end{complexbox}

\subsection*{ENSEMBLES DISJOINTS (Union-Find)}

\begin{defbox}
\textbf{Union-Find}
\begin{itemize}
\item Partition d'un ensemble
\item Sous-ensembles disjoints
\item Évolution : singletons → union
\end{itemize}
\end{defbox}

\begin{algobox}
\textbf{Opérations}
\begin{itemize}
\item \texttt{trouverClasse(e)}
\item \texttt{union(C1, C2)}
\end{itemize}
\end{algobox}

\begin{importbox}
\textbf{Implémentation 1 : Tableau}
\begin{itemize}
\item Trouver : $O(1)$
\item Union : $O(n)$
\end{itemize}

\textbf{Implémentation 2 : Forêt}
\begin{itemize}
\item Trouver : $O(h)$
\item Union : $O(h)$
\end{itemize}
\end{importbox}

\begin{defbox}
\textbf{Optimisations}
\begin{itemize}
\item \textbf{Union par rang} : attacher petit sous grand
\item \textbf{Compression de chemin} : pointer vers racine
\end{itemize}
\end{defbox}

\begin{complexbox}
\textbf{Avec optimisations}\\
$m$ opérations sur $n$ éléments :\\
$O(m \cdot \alpha(n))$ \\
$\alpha(n)$ = inverse Ackermann\\
$\approx O(1)$ amorti
\end{complexbox}

\vfill

\section*{\color{titrebleu}APPLICATIONS}

\begin{importbox}
\textbf{Pile} : parenthèses, undo\\
\textbf{File} : BFS, processus\\
\textbf{ABR} : dictionnaires triés\\
\textbf{Tas} : files de priorité, heapsort\\
\textbf{Union-Find} : Kruskal, cycles
\end{importbox}

\end{multicols}

\end{document}